\section{When Is a Public Hub Efficient?}
\label{sec:hub_efficiency}

This section isolates when a public hub is socially efficient \emph{as a stand-alone local policy}. Throughout, we set
\[
b=0.
\]
This makes it possible to evaluate the hub's local coordination value before introducing pipeline policy in Section~\ref{sec:refresh}.

\subsection{A Benchmark Efficiency Criterion}

Using the period-1 and period-2 local surplus expressions from Section~\ref{sec:model}, the net welfare gain from establishing a hub without pipeline policy is
\begin{equation}
\Delta_H
=
v_L P_A \bigl(m_A^1-m_A^0\bigr)
\Bigl[(1+\beta)-\beta\delta\bigl(m_A^1+m_A^0\bigr)\Bigr]
-F.
\label{eq:DeltaH_repeat}
\end{equation}

It is useful to define the hub's \emph{dynamic viability factor} as
\begin{equation}
\Lambda_A
:=
(1+\beta)-\beta\delta\bigl(m_A^1+m_A^0\bigr).
\label{eq:Lambda_def}
\end{equation}
Then \eqref{eq:DeltaH_repeat} becomes
\begin{equation}
\Delta_H
=
v_L P_A \bigl(m_A^1-m_A^0\bigr)\Lambda_A - F.
\label{eq:DeltaH_Lambda}
\end{equation}

The term $m_A^1-m_A^0$ captures the hub's immediate coordination gain, while $\Lambda_A$ summarizes whether that gain survives second-period redundancy. A larger $\delta$ lowers $\Lambda_A$ and therefore weakens the case for a hub.

The hub is socially efficient if and only if
\[
\Delta_H \ge 0.
\]

\begin{proposition}
\label{prop:hub_efficiency}
Consider the benchmark case without pipeline policy, $b=0$.

\begin{enumerate}[leftmargin=2em]
    \item If
    \[
    \Lambda_A \le 0,
    \]
    then the hub is never socially efficient:
    \[
    \Delta_H < 0
    \qquad
    \text{for all } N_A \ge 2 \text{ and } F>0.
    \]

    \item If
    \[
    \Lambda_A > 0,
    \]
    then the hub is socially efficient if and only if
    \begin{equation}
    P_A
    \ge
    \frac{F}{v_L\bigl(m_A^1-m_A^0\bigr)\Lambda_A}.
    \label{eq:pair_threshold}
    \end{equation}
    Equivalently, treating $N_A$ as continuous, the hub is efficient if and only if
    \begin{equation}
    N_A \ge N_H^*
    :=
    \frac{1}{2}
    \left(
    1+
    \sqrt{
    1+
    \frac{8F}{v_L\bigl(m_A^1-m_A^0\bigr)\Lambda_A}
    }
    \right).
    \label{eq:Nstar}
    \end{equation}
    If $N_A$ is restricted to integers, replace $N_H^*$ by its ceiling.
\end{enumerate}
\end{proposition}

Proposition \ref{prop:hub_efficiency} says that a hub creates value only if the gross local coordination benefit,
\[
v_L P_A \bigl(m_A^1-m_A^0\bigr),
\]
survives the dynamic penalty summarized by $\Lambda_A$ and is large enough to cover $F$. Failure can therefore come either from thin local scale or from strong dynamic redundancy.

An implication of the two-region setup is that the stand-alone criterion in \eqref{eq:DeltaH_Lambda} depends only on \emph{local} primitives. Outside-pool variables affect welfare levels, but they cancel out of the incremental comparison between a hub alone and no hub.

\subsection{Equivalent Threshold Characterizations}

Proposition \ref{prop:hub_efficiency} can be rewritten in two equivalent threshold forms.

First, for given $(N_A,m_A^0,m_A^1,\beta,\delta)$, define the critical fixed cost
\begin{equation}
F_H^*(N_A)
:=
v_L P_A \bigl(m_A^1-m_A^0\bigr)\Lambda_A.
\label{eq:Fstar}
\end{equation}
Then the hub is efficient if and only if
\[
F \le F_H^*(N_A).
\]

Second, for given $(N_A,m_A^0,m_A^1,\beta,F)$, define the critical redundancy rate
\begin{equation}
\delta_H^*(N_A)
:=
\frac{1+\beta}{\beta\bigl(m_A^1+m_A^0\bigr)}
-
\frac{F}{\beta v_L P_A\bigl(m_A^1-m_A^0\bigr)\bigl(m_A^1+m_A^0\bigr)}.
\label{eq:deltastar}
\end{equation}
Then the hub is efficient if and only if
\[
\delta \le \delta_H^*(N_A),
\]
provided the right-hand side is positive.

These thresholds show that stand-alone hub efficiency is governed by five local primitives: $N_A$, $m_A^1-m_A^0$, $v_L$, $F$, and $\delta$.

\subsection{Comparative Statics}

The next proposition summarizes the main comparative statics.

\begin{proposition}
\label{prop:comparative_statics_hub}
Suppose $\Lambda_A>0$. Then the hub-alone welfare gain $\Delta_H$ satisfies
\begin{align}
\frac{\partial \Delta_H}{\partial N_A}
&=
\frac{v_L}{2}(2N_A-1)\bigl(m_A^1-m_A^0\bigr)\Lambda_A
>0,
\label{eq:dDelta_dN}
\\[0.4em]
\frac{\partial \Delta_H}{\partial F}
&=
-1
<0,
\label{eq:dDelta_dF}
\\[0.4em]
\frac{\partial \Delta_H}{\partial \delta}
&=
-\beta v_L P_A \bigl(m_A^1-m_A^0\bigr)\bigl(m_A^1+m_A^0\bigr)
<0.
\label{eq:dDelta_ddelta}
\end{align}
Moreover, under the regularity condition \eqref{eq:regularity},
\begin{align}
\frac{\partial \Delta_H}{\partial m_A^1}
&=
v_L P_A\Bigl[(1+\beta)-2\beta\delta m_A^1\Bigr]
>0,
\label{eq:dDelta_dm1}
\\[0.4em]
\frac{\partial \Delta_H}{\partial m_A^0}
&=
- v_L P_A\Bigl[(1+\beta)-2\beta\delta m_A^0\Bigr]
<0.
\label{eq:dDelta_dm0}
\end{align}
\end{proposition}

Proposition \ref{prop:comparative_statics_hub} yields three immediate implications. First, the hub is more likely to be efficient in a thicker target ecosystem. Since
\[
P_A=\frac{N_A(N_A-1)}{2},
\]
the number of potential local pairs rises more than proportionally with $N_A$. Thin ecosystems are therefore disadvantaged even when the hub raises matching.

Second, faster local redundancy makes the hub less attractive. This is the model's reduced-form representation of dynamic redundancy risk from repeated interaction inside region $A$.

Third, a larger matching effect $m_A^1$ raises the hub's value, while a higher baseline $m_A^0$ lowers its incremental contribution. The hub is therefore most useful where pre-existing coordination is weak but the policy can materially reduce frictions.

\subsection{Discussion}

Two points matter for what follows. Hub efficiency is not determined by local scale alone, because strong redundancy can overturn even sizable short-run coordination gains. And the relevant object is $\Lambda_A$, not just the static increment $m_A^1-m_A^0$. A hub may therefore look successful on visible activity grounds while remaining dynamically weak. The next section compares that stand-alone logic with decentralized support under the model's dynamic coordination externality.
